\label{sec:experiments}

Fourteen experiments were conducted between 2026-06-28 and 2026-07-02 to map the reward and action-space design required for cloud-aware autonomous pointing.
Each run tests one concrete design choice; the sequence is cumulative rather than a search for a single winning configuration.
The arc divides into three phases: (i) diagnosing why naive training setups fail (Exps~1--5),
(ii) reward shaping and action-space design (Exps~4--9), and
(iii) MPO-specific stabilisation and further ablations (Exps~8--12), followed by scale and cadence probes (Exps~13--14).
Experiment~6 (modular encoder rerun) was deferred and is not included.

\paragraph{Experimental protocol.}
\label{sec:exp-protocol}
A fixed random seed (7) ensures identical altitude draws, cloud layouts, and target positions across all compared runs.
All experiments share the canonical simulation architecture described in Section~\ref{sec:arch}.

Before training begins, \emph{warmup episodes} pre-fill the replay buffer (Section~\ref{sec:replay-buffer}) using the deterministic baseline policy.
This was not the initial choice: early attempts using random torque commands or a zero-torque coast strategy yielded no reward signal, because neither strategy ever issues capture commands that earn positive credit.
The baseline policy, by contrast, reliably captures targets on clear passes.
Even so, only around ten capture actions fall within a 516-step episode ($\approx 2\%$ of all transitions), so the buffer is extremely sparse in rewarded transitions.
The warmup data therefore serves a dual purpose: it bootstraps the replay buffer with at least some positive-credit transitions, and the episode return it achieves is the KPI comparison anchor used throughout.

\emph{Evaluation episodes} run the current learned policy without gradient updates; eval return is the primary metric.
A run is said to show a \emph{learning signal} if its per-episode training return trends upward over training, indicating the policy has discovered the reward structure.
Concretely, this requires at least one training episode to substantially exceed the warmup return and for the trend to be sustained rather than a single outlier.
Runs without a learning signal converge to a flat return plateau from episode~1 onward, regardless of the number of training episodes, and are reported as \emph{not supported}.
In the verdict column, \emph{supported} means the design choice under test produced the expected learning or behavioural change (for example selective shuttering or a stable upward return trend); it does not by itself claim mission-ready outperformance of the baseline.
Results are archived as JSON files per experiment; key metrics are cited directly below.

% ─────────────────────────────────────────────────────────────────────────────
\subsection{Master comparison table}
\label{sec:exp-table}

Table~\ref{tab:exp-comparison} lists the full configuration and primary KPI for every run.
The \textbf{baseline anchor} row is the deterministic baseline evaluated at matched seeds (mean eval return $-38.6$; anchored in Exp~13).

\begin{table}[htbp]
  \centering
  \caption{Master experiment comparison. Algorithm: SAC = Soft Actor--Critic; MPO = Maximum a Posteriori Policy Optimisation.
    Action: T = torque (direct RW), V = vector pointing (nadir-relative offset).
    Reward flags: B = budget-exhausted shutter penalty; W = within-budget waste penalty; E = torque effort;
    S = safe-mode per-step penalty.
    Seeds: fixed 7 for all ML experiments; env draws are seeded per warmup fingerprint.
    Eval return is mean over 2 evaluation episodes unless otherwise noted.
    $\Delta$ is relative to the immediate predecessor configuration.}
  \label{tab:exp-comparison}
  \scriptsize
  \setlength{\tabcolsep}{4pt}
  \begin{tabular}{lllllccccc}
    \toprule
    Exp & Config & Alg & Act & Reward flags & Warmup & Train & Eval eps & Eval return & Verdict \\
    \midrule
    nb07      & baseline           & ---  & V  & ---        & ---  & ---  & 2 & $-38.6$   & --- \\
    \midrule
    1         & mpo\_t05           & MPO  & T  & ---        & 5    & 7    & 2 & $-101.6$  & not\_supported \\
    1         & mpo\_t09           & MPO  & T  & ---        & 5    & 7    & 2 & $-101.6$  & not\_supported \\
    2         & sac\_a0            & SAC  & T  & ---        & 5    & 7    & 2 & $-78.9$   & not\_supported \\
    2         & sac\_a1            & SAC  & T  & ---        & 5    & 7    & 2 & $-86.5$   & not\_supported \\
    3         & compare\_sac       & SAC  & T  & ---        & 5    & 37$^\dagger$  & 2 & $-22.2$   & supported \\
    3         & compare\_mpo       & MPO  & T  & dense      & 5    & 22$^\dagger$  & 2 & $-51.6$   & not\_supported \\
    4         & ref0 (torque)      & SAC  & T  & ---        & 5    & 50   & 2 & $-11.1$   & partial \\
    4         & ref1 (vector)      & SAC  & V  & ---        & 5    & 50   & 2 & $-24.1$   & partial \\
    5         & mpo\_s (90/140)    & MPO  & T  & sparse     & 5    & 50   & 2 & $-51.6$   & not\_supported \\
    5         & mpo\_m (140/256)   & MPO  & T  & sparse     & 5    & 50   & 2 & $-51.6$   & not\_supported \\
    5         & mpo\_l (256/512)   & MPO  & T  & sparse     & 5    & 50   & 2 & $-51.6$   & not\_supported \\
    7         & penalty\_on        & SAC  & V  & B          & 5    & 50   & 2 & $+10.6$   & supported \\
    8         & torque\_sparse     & MPO* & T  & ---        & 5    & 50   & 2 & $-81.1$   & partial \\
    9         & waste\_off\_budget\_on & SAC & V & B (W off)& 5    & 50   & 2 & $+106.4$  & supported \\
    10        & safe\_mode\_on     & MPO* & T  & S($k$=2.0) & 5    & 50   & 2 & $-510.4$  & not\_supported \\
    11        & vector\_sparse     & MPO* & V  & ---        & 5    & 50   & 2 & $+45.5$   & supported \\
    12        & vector\_effort     & MPO* & V  & E($k$=0.1) & 5    & 50   & 2 & $+25.3$   & not\_supported \\
    13        & baseline\_1\_1\_1  & nb07 & V  & ---        & ---  & ---  & 2 & $-38.6$   & --- \\
    13        & cadence\_1\_1\_10  & MPO* & V  & ---        & 5    & 50   & 2 & $-46.4$   & not\_supported \\
    13        & cadence\_1\_1\_50  & MPO* & V  & ---        & 5    & 50   & 2 & $-15.4$   & partial \\
    14        & multienv           & MPO* & V  & ---        & 5    & 350  & 2 & n/a$^\ddagger$ & not\_supported \\
    \bottomrule
    \multicolumn{10}{l}{$^\dagger$ Early abort (patience 20 eps). \quad *MPO with decoupled-KL dual fix (Exp~8). \quad $^\ddagger$ Run crashed (OSError).} \\
  \end{tabular}
\end{table}

% ─────────────────────────────────────────────────────────────────────────────
\subsection{Exp~1: Shutter threshold ablation}
\label{sec:exp1}

\paragraph{Motivation.}
Preliminary runs showed the MPO agent emitting $\sim$968 shutter commands per episode, far exceeding the budget of ten and firing the shutter on nearly every controller step.
The hypothesis was that raising the continuous action threshold (above which a shutter command is issued) from 0.5 to 0.9 would mechanically reduce this saturation
and, combined with a 15\,s credit window around each accepted capture, create a learnable reward gradient.

\paragraph{Config and result.}
Two configurations (shutter threshold 0.5 and threshold 0.9) were run for 7 training episodes at $\Delta t=1.5$\,s.
Both produced an identical eval return of $-101.6$ with $\sim$516 shutter commands per episode.
The policy's shutter action was saturated near its maximum value (mean $\approx 0.999$) at both thresholds.
It fired the shutter on almost every step regardless of threshold setting.
Neither showed a learning signal; every shutter command produced zero useful capture credit.

\paragraph{Verdict.}
\textbf{Not supported.}
Threshold tuning is not an effective lever while the policy saturates the action dimension.
The credit window creates latent mid-episode reward bursts but does not fix credit assignment.

% ─────────────────────────────────────────────────────────────────────────────
\subsection{Exp~2: Modular encoder ablation}
\label{sec:exp2}

\paragraph{Motivation.}
The controller observation is 105-dimensional with 50+50 structured bearing/mask fields.
Flattening these into a single MLP may be sample-inefficient compared to a group-compressor
\citep{gorishniy2022revisiting}.
Variant A1 replaced the flat concat with a $\mathrm{Linear}(50 \to 8)$ compressor on each field group.

\paragraph{Config and result.}
SAC sparse, 7 training episodes.
A0 (flat): eval $-78.9$, 223 shutters/ep.
A1 (compress): eval $-86.5$, 98 shutters/ep.
Neither variant showed a learning signal; both produced zero meaningful captures.
A1's critic loss diverged (3190 vs A0's 1420 at episode~7), indicating training instability in the compressed branch.

\paragraph{Verdict.}
\textbf{Not supported.}
The structured observation encoder shows no benefit in the 7-episode non-learnable slice.
The encoder is deferred for a rerun once a stable learning configuration exists (deferred Exp~6).

% ─────────────────────────────────────────────────────────────────────────────
\subsection{Exp~3: SAC vs.\ MPO algorithm comparison}
\label{sec:exp3}

\paragraph{Motivation.}
Exps~1--2 both used MPO; learning failure could be algorithm-specific.
Exp~3 runs SAC (sparse reward) and MPO (dense reward) under matched protocol to isolate the
algorithm--reward interaction \citep{haarnoja2018sac,abdolmaleki2018mpo}.

\paragraph{Config and result.}
Early-abort patience of 20 episodes.
SAC sparse (early-abort at ep~37): eval $-22.2$, clear learning signal, best training return $+5.3$.
MPO dense (early-abort at ep~22): eval $-51.6$, no learning signal, KL divergence $\approx 3.8\times10^5$.
SAC clearly separated from MPO: the SAC policy converged while MPO's policy update blew up.

\paragraph{Verdict.}
\textbf{Supported.}
SAC learns under sparse reward; MPO dense fails in this configuration.
Algorithm choice matters; MPO's failure traced to an unconstrained M-step dual, motivating the fix in Exp~8.

% ─────────────────────────────────────────────────────────────────────────────
\subsection{Exp~4: Vector pointing mode}
\label{sec:exp4}

\paragraph{Motivation.}
Outputting a direct reaction-wheel torque conflates two problems: target scheduling and rate control.
The vector pointing interface delegates rate control to a proportional--derivative loop
and exposes a nadir-relative offset to the policy, reducing the action space the learner must shape for scheduling.

\paragraph{Config and result.}
SAC sparse, 50 training episodes.
Ref0 (torque): eval $-11.1$, clear learning signal, best training return $+8.4$.
Ref1 (vector): eval $-24.1$, clear learning signal, best training return $+95.8$.
Ref1's eval was worse, but video of training episode~11 (+95.8) showed deliberate off-nadir scheduling
and sparse mid-orbit shutter firings.
This showed a qualitatively different and promising behaviour.
End-of-episode shutter spam after budget exhaustion remained apparent in Ref1.

\paragraph{Verdict.}
\textbf{Partial.}
Vector mode is valid and both variants learn, but the eval gap between them is unexplained.
The shutter-budget ignorance motivates Exp~7.
Vector pointing mode is retained as the preferred action interface for subsequent experiments.

% ─────────────────────────────────────────────────────────────────────────────
\subsection{Exp~5: MPO network width ablation}
\label{sec:exp5}

\paragraph{Motivation.}
Before investigating the MPO dual, the hypothesis that a 90/140-unit head was capacity-limited
was tested with three widths: S (90/140), M (140/256), L (256/512).
This was gated after the encoder experiment following \citet{bruin2018irl} suggestions
on structure-before-capacity.

\paragraph{Config and result.}
MPO sparse, 50 training episodes.
All three widths produced identical eval $-51.6$ and identical KL freeze/explosion patterns.
Post-episode-1 returns were flat; every shutter command produced zero useful capture credit across all three width variants.

\paragraph{Verdict.}
\textbf{Not supported.}
Network width is not the bottleneck; the shared collapse pattern points to the learning algorithm.

% ─────────────────────────────────────────────────────────────────────────────
\subsection{Exp~7: SAC vector + budget-exhausted shutter penalty}
\label{sec:exp7}

\paragraph{Motivation.}
Exp~4 video (Ref1 ep~11) showed mid-orbit sparse shuttering but persistent shutter spam after the
10-capture budget was exhausted.
The budget-exhausted penalty adds $-k_\text{waste}$ to the step reward whenever the agent commands a
shutter with zero remaining budget (Section~\ref{sec:reward}).
Torque-path SAC with a static threshold (Exp~1) failed to stop spam; a direct negative reward on
the offending action is the minimal correct lever.

\paragraph{Config and result.}
SAC sparse + vector + budget penalty, 50 training episodes.
Eval $+10.6$ (vs Ref1 baseline $-24.1$, $\Delta=+34.7$).
Train shutter commands: 478/ep vs 7071/ep Ref1 ($-93\%$).
Train meaningful shutter fraction: 0.220 vs 0.015 Ref1 ($15\times$).
Post-peak return std: 17.5 vs 28.4 ($-38\%$).

\paragraph{Verdict.}
\textbf{Supported.}
Budget-exhausted penalty is an effective design lever: shutter spam drops sharply and selective scheduling emerges.
This reward term is kept for all subsequent experiments.

% ─────────────────────────────────────────────────────────────────────────────
\subsection{Exp~8: MPO decoupled-KL dual fix (torque mode)}
\label{sec:exp8}

\paragraph{Motivation.}
Exps~3 and~5 both showed MPO freezing or exploding at KL $\approx 3.8\times10^5$ or $\approx 0$.
Investigation identified the root cause: the M-step used a single coupled Lagrange multiplier
for both mean and variance KL constraints, instead of the two independent multipliers specified
in \citet{abdolmaleki2018mpo}.
Exp~8 validates the corrected decoupled-KL dual in torque mode against the Exps~3/5 comparators.

\paragraph{Config and result.}
MPO (fixed dual) sparse torque, 50 training episodes, seed~7.
Clear learning signal; eval $-81.1$ vs pre-fix floor $-51.6$.
Best train $+8.7$ (ep~28).
KL converged to 0.021; temperature dual settled at 0.788.
Train torque saturation fraction: 0.951 (vs $\approx 0.998$ pre-fix).
No meaningful eval captures.

\paragraph{Verdict.}
\textbf{Partial.}
The algorithm fix is confirmed: bounded KL and rising train returns.
Eval return remains below the baseline anchor; torque mode still saturates and safe-mode activates frequently, motivating Exps~10 and~11.

% ─────────────────────────────────────────────────────────────────────────────
\subsection{Exp~9: SAC shutter reward split (waste penalty off)}
\label{sec:exp9}

\paragraph{Motivation.}
Exp~7 had both the budget-exhausted penalty and the within-budget waste penalty active.
Since applied-capture credit already scales with image quality and cloud fraction (Eq.~\ref{eq:r-cap}),
a second explicit waste penalty on the same action may over-constrain shutter exploration.
Exp~9 tests turning off the waste penalty while keeping the budget-exhausted term.

\paragraph{Config and result.}
SAC sparse + vector + budget penalty ($W$ off), 50 training episodes.
Eval $+106.4$ (vs Exp~7 $+10.6$, $\Delta=+95.8$).
Best train $+169.9$ (ep~36).
Train meaningful fraction: 0.286 vs 0.220 (Exp~7).
Eval meaningful fraction: 0.667 vs 0.400.
Post-budget shutter commands: 75 (penalty still active).

\paragraph{Verdict.}
\textbf{Supported.}
Relaxing the over-constraining within-budget waste penalty yields a large eval improvement.
This is the best SAC configuration in the series and defines the reward stack used as the reference for later MPO comparisons.

% ─────────────────────────────────────────────────────────────────────────────
\subsection{Exp~10: MPO safe-mode penalty}
\label{sec:exp10}

\paragraph{Motivation.}
Exp~8 showed safe-mode activating 5--6 times per episode in torque mode with no reward signal.
A per-step penalty during safe-mode recovery intervals ($k=2.0$) was hypothesised to reduce activations
and improve applied torque credit assignment.

\paragraph{Config and result.}
MPO (fixed dual) sparse torque + safe-mode penalty ($k=2.0$), 50 training episodes.
Eval $-510.4$ (vs Exp~8 $-81.1$, catastrophic regression).
Training return improved from $-1168$ (ep~0) to $\approx-350$ (ep~23) but stagnated.
Train shutter commands dropped from 153 (ep~0) to 2--7 (eps~30--49).

\paragraph{Verdict.}
\textbf{Not supported.}
The penalty taught the agent to avoid all aggressive torque (including captures) rather than
only avoiding safe-mode entry.
Safe-mode activations co-occur with off-nadir pointing required for capture; the penalty suppressed both.
Per-step safe-mode penalty at this coefficient is ruled out for torque-mode MPO.
Vector mode (Exp~11) provides a structurally cleaner solution.

% ─────────────────────────────────────────────────────────────────────────────
\subsection{Exp~11: MPO fixed dual, vector mode}
\label{sec:exp11}

\paragraph{Motivation.}
Exp~10 ruled out reward-based safe-mode mitigation in torque mode.
Exp~8 confirmed the dual fix; the natural next step is MPO in vector pointing mode,
mirroring the Exp~4 SAC comparison but now with a working MPO agent.
Vector mode avoids the safe-mode interference seen in torque mode.

\paragraph{Config and result.}
MPO (fixed dual) sparse vector, 50 training episodes, seed~7.
Eval $+45.5$ (vs Exp~8 $-81.1$, $\Delta=+126.6$).
Best train $+82.3$ (ep~42).
Positive returns from episode~7 onward; monotonic improvement from ep~5.
KL: 0.014; $\eta$: 0.014 (vs 0.788 in torque mode).
Shutter commands: 148 (ep~0) $\to$ 8--10 (eps~30--49).
The policy learned selective shuttering.

\paragraph{Verdict.}
\textbf{Supported.}
Vector mode resolves the torque-mode MPO stagnation: the first MPO configuration with positive eval return.
The temperature dual settles much lower in vector mode, reflecting a tighter value distribution in the simplified action space.
Vector pointing mode is retained for subsequent MPO experiments.
The remaining gap to SAC Exp~9 (+106.4) suggests the MPO branch has not yet been tested with the same reward stack as Exp~9.

% ─────────────────────────────────────────────────────────────────────────────
\subsection{Exp~12: MPO vector + torque-effort penalty ablation}
\label{sec:exp12}

\paragraph{Motivation.}
The training workflow disables the torque-effort regulariser in vector mode by default.
Exp~12 re-enables effort ($k=0.1$) while keeping all other Exp~11 settings.

\paragraph{Config and result.}
MPO (fixed dual) sparse vector + torque effort ($k=0.1$), 50 training episodes.
Eval $+25.3$ (vs Exp~11 $+45.5$, $\Delta=-20.2$).
Best train $+75.6$ (ep~20).
Four consecutive zero-return episodes at eps~36, 41, 42, 43.
KL: 0.013; $\eta$: 0.027.

\paragraph{Verdict.}
\textbf{Not supported.}
Torque effort in vector mode hurts performance: the effort penalty targets the pointing-controller wheel command,
not the policy's pointing offset, creating a conflicting gradient that discourages off-nadir pointing
irrespective of imaging context.
Effort regularisation should remain disabled in vector mode.

% ─────────────────────────────────────────────────────────────────────────────
\subsection{Exp~13: MPO learning-cadence sweep}
\label{sec:exp13}

\paragraph{Motivation.}
The MPO agent performs one gradient update per environment step by default.
If the agent learns faster than the environment provides signal, or if gradient noise from a partially
filled replay buffer impedes early learning, a different cadence may help.
Exp~13 evaluates three cadence profiles: update every step, every 10 steps, and every 50 steps.

\paragraph{Config and result (partial).}
MPO (fixed dual) vector sparse, 50 training episodes per configuration, 2 eval episodes.
The deterministic baseline (comparison anchor): eval $-38.6$.
Update every 10 steps: eval $-46.4$ (worse than baseline).
Update every 50 steps: eval $-15.4$ (best in sweep; still below Exp~11 $+45.5$).

\paragraph{Verdict.}
\textbf{Partial / incomplete.}
Sparser gradient updates (every 50 steps) marginally outperform the per-step default in this short slice,
but results remain below Exp~11 and the every-10-steps configuration falls below the comparison anchor ($-38.6$).
The full cadence sweep was still in progress at report cutoff; a longer training slice is needed for a definitive conclusion.

% ─────────────────────────────────────────────────────────────────────────────
\subsection{Exp~14: MPO multi-environment target selection}
\label{sec:exp14}

\paragraph{Motivation.}
All prior experiments use a single environment seed (7) with a fixed target corridor.
Generalisation to diverse orbital configurations, altitudes, and cloud draws requires
training across multiple environment draws.
Exp~14 attempted to run 350-episode multi-environment MPO to test this scaling hypothesis.

\paragraph{Config and result.}
MPO (fixed dual) vector sparse, 350 training episodes, multi-seed draws.
The run crashed during Stage~C with a filesystem error
(likely a path-length limit on Windows for long run-directory names).
Stage~B partial training metrics were available (mean score 0.228) but are not directly comparable
to the episodic eval return metric used in the rest of the series.

\paragraph{Verdict.}
\textbf{Not supported} (run failure; results inconclusive).
Multi-environment training remains an open direction.
Recommended follow-up: shorten run directory names, increase episode budget, and match the evaluation
metric to the rest of the series.
